\section{Methods}

\subsection{Model}\label{subsec:model}

Consider two interconnected regions with populations $N_1$ and $N_2$.
Let $S_i(t)$, $I_i(t)$, $R_i(t)$ denote the number of susceptible,
infected, and recovered individuals in region $i$ at time $t$, with
$S_i + I_i + R_i = N_i$. Regional coupling is governed by the contact
matrix
%
%
$$
\mathbf{C} =
\begin{bmatrix}
  1 - \theta  &     \theta \\
      \theta  & 1 - \theta
\end{bmatrix}
$$
%
%
where $\theta \in [0, 0.5]$ is the connectivity parameter. The
transmission rate is a standard seasonal forcing ($\delta$ is forcing
amplitude, $\varphi_i$ is a phase offset):
%
%
$$
\beta_i(t) = \beta_0(1 + \delta\sin(2\pi t + \varphi_i)),
$$
%
%
where $t$ is measured in years. Then, following~\citet{edlund2011,
  keeling2008} the force of infection in region $i$ is:
\begin{equation}\label{eq:foi}
\lambda_i(t) = \beta_i(t) \frac{\sum_j C_{ij}\, I_j(t)}{\sum_j C_{ij}\, N_j}
\end{equation}

The expected incidence is $\mu_i(t) = S_i(t)[1 -
  \exp(-\lambda_i(t))]$, and the recovery probability per time step is
$\exp(-\gamma)$. Time steps are denoted $\Delta t$, hence the
discrete-time SIR dynamics are:
\begin{align}\begin{split}\label{eq:discrete_sir}
S_i(t+\Delta t) &= S_i(t) \exp(-\lambda_i(t)) \\
I_i(t+ \Delta t) &= I_i(t) \exp(-\gamma) + \mu_i(t)
\end{split}\end{align}
We use SIR rather than SIRS dynamics because we conduct inference for
each influenza season separately. Model parameters are calibrated to
seasonal influenza following~\citet{shaman2010}; see
Table~\ref{tab:params}.

\begin{table}[h]
\centering
\begin{tabular}{llll}
\toprule
Parameter & Value & Description & Source or calculation\\
\midrule
$R^{\text{max}}_0$ & 3.52 & $R_0$ at peak season & \cite{shaman2010} \\
$R^{\text{min}}_0$ & 1.12 & $R_0$ during the summer & \cite{shaman2010} \\
$D$ & 3.24 & Mean infection period (days) & \cite{shaman2010} \\
$\gamma$ & $2.16$ & Weekly recovery rate & $7 / D$ \\
$\beta_{\text{max}}$ & $3.11$ & Peak-season transmission & $R_0^{\text{max}}(1 - e^{-\gamma})$\\
$\beta_{\text{min}}$ & $0.99$ & Low-season transmission & $R_0^{\text{min}}(1 - e^{-\gamma})$ \\
$\beta_0$ & $2.05$ & Mean transmission rate & $(\beta_{\text{min}} + \beta_{\text{max}})/2$ \\
$\delta$ & $0.517$ & Seasonal amplitude & $\beta_{\text{max}} / \beta_0 - 1$\\
$\rho$ & $5.36 \times 10^{-4}$ & Infection Fatality Rate & \citet{rolfes2018}\\
$\varphi$ & & Phase of seasonal forcing & per state \\
\bottomrule
\end{tabular}
\caption{Model parameters calibrated to seasonal influenza with
  $\beta(t) = \beta_0(1+\delta \sin(2\pi t + \varphi))$.}
\label{tab:params}
\end{table}

\subsection{Observation Model and Likelihood}\label{subsec:noise}
We assume each infection is reported independently with probability
$\rho$, giving rise to binomially distributed observations, which we approximate as Gaussian:
%
%
$$
Y_i(t) = \rho\mu_i(t) + \epsilon_i(t)
$$
%
%
with
%
%
$$
\epsilon_i(t) \sim \mathcal{N}(0, \rho(1{-}\rho)\mu_i(t)).
$$
%
%
%% The Gaussian model is a deliberate simplification. Real
%% surveillance data typically exhibit overdispersion; a negative
%% binomial model with variance $\rho\mu_i + (\rho\mu_i)^2/k$ would be
%% more realistic. At any fixed parameter value, overdispersion
%% increases the observation variance, which reduces the FIM weights
%% and inflates the CRLB. While the MLE itself may shift under a
%% different noise model, the qualitative conclusion---that the CRLB
%% for $\theta$ is large when epidemics are synchronized---is robust,
%% since the underlying mechanism (vanishing sensitivity
%% $\partial\mu/\partial\theta$ under symmetry) is independent of the
%% observation model.

The unknown parameters for a single season are $\boldsymbol{\phi} =
(\theta, S_1(0), I_1(0), S_2(0), I_2(0))$, giving the log-likelihood:
\begin{equation}\label{eq:loglik}
\ell(\boldsymbol{\phi}) = -\frac{1}{2}\sum_{t=0}^{T-1}\sum_{i=1}^{2}
\left[\log(2\pi \sigma_i^2(t)) + \frac{(Y_i(t) -
    \rho\mu_i(t))^2}{\sigma_i^2(t)}\right]
\end{equation}
where $\sigma_i^2(t) = \rho(1{-}\rho)\mu_i(t)$.

\subsection{Derivatives}\label{subsec:sens}
Computing the CRLB and optimizing the likelihood both require
partial derivatives of model outputs with respect to parameters. Let $\phi \in
\{\theta, S_1(0), S_2(0), I_1(0), I_2(0)\}$ denote a parameter ($\phi
= \phi_k$ for $k = 1, \dots, 5$). The partial derivative of the force
of infection~\eqref{eq:foi} is:
\begin{equation}\label{eq:sens_foi}
\frac{\partial \lambda_i}{\partial \phi} = \beta_i(t) \left[
\frac{1}{(\mathbf{C}\mathbf{N})_i}\left(
  \frac{\partial \mathbf{C}}{\partial \phi} \mathbf{I}
  + \mathbf{C} \frac{\partial \mathbf{I}}{\partial \phi}\right)_i
- \frac{\lambda_i}{\beta_i (\mathbf{C}\mathbf{N})_i}
  \left(\frac{\partial \mathbf{C}}{\partial \phi} \mathbf{N}\right)_i
\right]
\end{equation}
The contact matrix derivative is

\begin{align*}
\partial \mathbf{C}/\partial \theta =
\begin{bmatrix}
  -1 &  1 \\
  \phantom{-}1 & -1
\end{bmatrix},\\
\qquad
\partial \mathbf{C}/\partial \phi = \mathbf{0} \text{ for } \phi \neq \theta,
\end{align*}
%% so for $\phi\neq\theta$ only the
%% $\mathbf{C}\,\partial\mathbf{I}/\partial\phi$ term in
%% eq.~\eqref{eq:sens_foi} survives.
Partial derivatives of the expected incidence are found from the
product rule:
\begin{equation}\label{eq:sens_mu}
  \frac{\partial \mu_i}{\partial \phi} =
  %
  %
  \frac{\partial S_i}{\partial \phi} [1 - e^{-\lambda_i}] + S_i
  e^{-\lambda_i} \frac{\partial \lambda_i}{\partial \phi}.
\end{equation}
Partials of the system state propagate via:
\begin{align}\begin{split}\label{eq:sens}
\frac{\partial S_i(t+\Delta t)}{\partial \phi} &= e^{-\lambda_i}
\frac{\partial S_i}{\partial \phi} - S_i e^{-\lambda_i}
\frac{\partial \lambda_i}{\partial \phi} \\
\frac{\partial I_i(t+\Delta t)}{\partial \phi} &= e^{-\gamma}
\frac{\partial I_i}{\partial \phi} +
\frac{\partial \mu_i}{\partial \phi}
\end{split}\end{align}
with initial conditions $\partial S_i(0)/\partial S_j(0) =
\delta_{ij}$, $\partial I_i(0)/\partial I_j(0) = \delta_{ij}$, and
zero otherwise. We find these partial derivatives jointly with the
state equations.

\subsection{Cram\'er-Rao Lower Bound}\label{subsec:crlb}
For a single epidemic season, the Cram\'er-Rao Lower Bound (CRLB)
states that $\text{Var}(\hat\phi_k) \geq [\mathbf{J}^{-1}]_{kk}$,
where $\mathbf{J}$ is the $5 \times 5$ Fisher Information Matrix with
entries $J_{kl} =
\mathbb{E}[(\partial\ell/\partial\phi_k)(\partial\ell/\partial\phi_l)]$.
One can verify that if we define
\begin{equation*}
A_i(t) := -\frac{1}{2\mu_i(t)} + \frac{Y_i(t) -
  \rho\mu_i(t)}{(1{-}\rho)\mu_i(t)} + \frac{(Y_i(t) -
  \rho\mu_i(t))^2}{2\rho(1{-}\rho)\mu_i(t)^2},
\end{equation*}
then differentiating the log-likelihood~\eqref{eq:loglik} yields
%
%
\begin{equation}\label{eq:grad}
\partial \ell / \partial \phi = \sum_{t,i} A_i(t)\,\partial
\mu_i(t)/\partial \phi.
\end{equation}
%
%
We can also easily verify $\mathbb{E}[A_i(t)] = 0$. Since observations
are independent across $(i,t)$, we obtain, using moments of the
Gaussian $Y_i(t) -\rho \mu_i(t)$ and eq.~\eqref{eq:grad}:
\begin{align}\label{eq:fim_scalar}
  \begin{split}
    J_{kl} &= \mathbb{E}\left [\frac{\partial \ell}{\partial \phi_k}\frac{\partial \ell}{\partial \phi_l}\right ] \\
    %
    %
    %
    &=\sum_{t,i} \mathbb{E}[A_i^2(t)] \frac{\partial \mu_i(t)}{\partial \phi_k}  \frac{\partial \mu_i(t)}{\partial \phi_l} \\
    %
    %
    %
    &= \frac{1+\rho}{2(1-\rho)}\sum_{t=0}^{T-1}\sum_{i=1}^{2}
    \frac{1}{\mu_i(t)^2} \frac{\partial \mu_i(t)}{\partial \phi_k}
    \frac{\partial \mu_i(t)}{\partial \phi_l}
  \end{split}
\end{align}
%
%
In matrix form, eq.~\eqref{eq:fim_scalar} can be written as 
\begin{align}\label{eq:fim}
  \begin{split}
  \mathbf{J} &= \mathbf{G}^T \mathbf{W} \mathbf{G} \\
  &\text{ where } \\
  %
  %
  G_{(i,t),k} &= \partial \mu_i(t)/\partial \phi_k \\
  %
  %
  \mathbf{W} &= \frac{1+\rho}{2(1-\rho)}\,\text{diag}(\mu_1(0)^{-2},
  \ldots, \mu_2(T{-}1)^{-2}).
  \end{split}
\end{align}

When multiple seasons are available, we assume independence across
seasons, treating each season's initial conditions as nuisance
parameters. We aggregate to find the Cram\'er-Rao Lower Bound:
\begin{equation}\label{eq:crlb_multi}
\text{Var}(\hat{\theta}) \geq \left(\sum_{s} \frac{1}{
  [\mathbf{J}^{-1}(s)]_{11}}\right)^{-1},
\end{equation}
where $\mathbf{J}(s)$ is the Fisher information matrix for season $s$.
Seasons where the CRLB computation fails (e.g., due to near-zero
incidence producing a singular Fisher information matrix) contribute
no information and are assigned $[\mathbf{J}^{-1}(s)]_{11} = \infty$.

Under standard regularity conditions, the MLE is asymptotically
efficient: $\hat\theta_{\text{MLE}} \sim \mathcal{N}(\theta,
[\mathbf{J}^{-1}]_{11})$ as the number of observations
grows~\citep{casella2024statistical}. The CRLB therefore also serves
as a practical approximation for the variance of the
maximum-likelihood estimate of $\theta$.

Identifiability in systems-biology models is sometimes explored via
the profile likelihood~\citep{raue2009} which is defined, for a
single season via
%
%
$$
\ell_{\text{PL}}(\theta) = \max_{S_1(0), S_2(0), I_1(0), I_2(0)} \ell(\theta, S_1(0), S_2(0), I_1(0), I_2(0)).
$$
%
%
The profile likelihood is calculated numerically over a grid of
$\theta$ values, optimizing over all other parameters at each grid
point. Such a calculation would be computationally prohibitive at
the scale of our 1081-pair analysis and the
$(\theta, \delta, \varphi, \text{IC})$ grid sweep used
in~\S\ref{subsec:sim_crlb}. The CRLB is effectively a quadratic
approximation to the profile likelihood around the maximum-likelihood
estimate; it is cheaper to compute, at the cost of losing the full
nonlinear shape that a profile likelihood would reveal in strongly
non-quadratic regimes.

\subsection{Maximum-Likelihood Estimation}\label{subsec:mle}
We maximize the log-likelihood using \texttt{SLSQP} from
\texttt{NLopt}). Analytical gradients were calculated via
eq.~\eqref{eq:grad}. To handle potential local minima, we use a
multi-start strategy with 200 random initial points and retain the
solution with the highest likelihood. Pilot experiments show that 20
initial points suffice, so we take an order of magniture more just to
be safe.

\subsection{Influenza Mortality Data}\label{subsec:real_data}
We analyze pneumonia and influenza (P\&I) mortality across 47 US
states.  We exclude Alaska and Hawaii since we do not have absolute
humidity data for them, and we exclude Arizona since it is a very arid
state with relatively inconsistent absolute humidity values (see
Supporting Information). We integrate CDC FluView data
(2010--2018)~\citep{cdcfluview2024} and HealthData.gov data
(2017--2025)~\citep{healthdatagov2025}, excluding COVID pandemic years
2020--2022. Alaska and Hawaii are excluded because they lack absolute
humidity data for phase estimation; Arizona is excluded because its
arid climate produces humidity patterns inconsistent with the
sinusoidal forcing model. Since P\&I deaths include a year-round
non-influenza baseline, we estimate influenza burden as excess
mortality above a pre-season baseline, computed as the mean weekly
death count during July--October preceding each flu
season~\citep{serfling1963, viboud2006}. The infection fatality rate
$\rho \approx 5.36 \times 10^{-4}$ follows from the first row of Table
1 in~\citet{rolfes2018}. Seasonal phases $\varphi_i$ are estimated
from absolute humidity data~\citep{daon2025, shaman2009} by maximizing
the correlation between $-\text{AH}(t)$ and $\sin(2\pi t +
\varphi)$. The median phase difference between state pairs is
approximately $\pi/45$, reflecting the near-perfect entrainment of
temperate US influenza seasons to shared seasonal forcing. We define
each flu season as $T = 20$ weekly observations starting
November~1$^{\text{st}}$ and fit a single connectivity $\theta$ shared
across all seasons, with per-season initial conditions $(S_i(0),
I_i(0))$ inferred via maximum-likelihood. We aggregate the per-season
Cram\'er-Rao lower bounds via eq.~\eqref{eq:crlb_multi} to obtain a
bound on the overall estimator variance for the connectivity $\theta$.

\subsection{Simulated CRLB Analysis}\label{subsec:sim_crlb}
To verify that the lack of identifiability observed for the real
data in~\S\ref{subsec:real_data} is driven by entrainment, we conduct
a simulated CRLB analysis using realistic initial conditions. We
compute the CRLB across a grid of connectivity values $\theta \in
[10^{-4}, 10^{-1}]$ and seasonal forcing amplitudes $\delta \in [0,
1)$, comparing synchronized ($\varphi_1 = \varphi_2$) and asynchronous
($\varphi_1 - \varphi_2 = \pi$) epidemics. At each grid point, the
initial conditions $(S_i(0)/N_i, I_i(0)/N_i)$ for each simulated
season and region are drawn independently from the empirical
distribution of per-season, per-state initial conditions fitted
in~\S\ref{subsec:real_data} (see SI~Fig.~\ref{fig:si_ics}). We repeat
each configuration many times with fresh IC draws and report the
median CRLB ratio (synchronized / asynchronous) as a function of
$(\theta, \delta)$.

\subsection{Linear Time-Series Benchmark: Univariate AR versus Bivariate VAR}\label{subsec:varbench}

To test the paper's central claim in a model class completely
distinct from the mechanistic SIR, we run a forecast-skill benchmark
that compares univariate autoregressive (AR) models to bivariate
vector autoregressive (VAR) models on weekly pneumonia and influenza
mortality. If inter-state \emph{directed predictive} dependence is
present in the data, it should manifest as improved out-of-sample
forecast skill for the bivariate model relative to two independent
univariate models.

\paragraph{Data.} We use continuous year-round raw P\&I weekly death
counts from~\citet{cdcfluview2024, healthdatagov2025} for the ten
largest US states by 2020 population (California, Texas, Florida,
New~York, Pennsylvania, Illinois, Ohio, Georgia, North~Carolina,
Michigan). Counts are normalised to per-100\,000-population weekly
rates by linearly interpolating the annual July-1 Census Bureau
population estimates, and we retain only weeks where all ten states
have non-missing data.

\paragraph{Models.} For each of the $\binom{10}{2}=45$ state pairs
$(y_1, y_2)$ and each candidate lag order $p$, we compare:
\begin{itemize}
\item \emph{Two independent univariate AR($p$) models}, each fit to
  one state's series, with intercept, linear trend, and 51 weekly
  seasonal dummies.
\item \emph{A single bivariate VAR($p$) model} fit to the pair, with
  the same intercept, trend, and seasonal-dummy structure on each
  equation. The bivariate model differs from the two univariates only
  by including cross-lag coefficients $a^{(k)}_{ij}$ ($i \neq j$,
  $k=1,\ldots,p$) and estimating a full $2\times 2$ residual
  covariance $\Sigma$.
\end{itemize}
Both models are fit by ordinary least squares on the same engine, so
their point forecasts use identical optimisation machinery. Lag
order $p$ is selected per pair by minimising AIC over
$p\in\{0,\ldots,8\}$ on the training split: independently for the two
univariates ($p_1, p_2$) and jointly for the bivariate ($p_\text{VAR}$).

\paragraph{Training and test splits.} We train on the decade 2010-2019
(522 weekly observations per state: continuous and pre-pandemic) and
test on the three-year window 2023-2025 (156 weekly observations),
skipping the COVID-19 pandemic period 2020-2022 consistent with the
main mechanistic analysis. For each test-period observation
$y_{i,t}$, we compute a one-step-ahead forecast~$\hat y_{i,t}$ using
the coefficients fitted on 2010-2019 and the actual lags
$y_{i,t-1},\ldots,y_{i,t-p}$ from the observed series; no test data
leak into coefficient estimation. We report, for each pair, the
training-set AIC difference $\Delta\text{AIC} = \text{AIC}_\text{VAR}
- (\text{AIC}_1 + \text{AIC}_2)$ and the test-set one-step-ahead
RMSE for each model pooled across both states and all test weeks.

\paragraph{What the comparison isolates.} A reduced-form VAR's point
forecast for each series uses only \emph{past} lags of both series;
the contemporaneous residual covariance $\Sigma$ does not enter the
forecast equation. Out-of-sample RMSE therefore isolates the
contribution of the cross-lag coefficients $a^{(k)}_{ij}$ ($i\neq j$)
-- whether state~A's past helps predict state~B's future beyond
state~B's own past. This is the AR-framework analogue of the
directed coupling parameter $\theta$ in the mechanistic SIR model.
By contrast, AIC penalises extra parameters but credits both
cross-lag terms \emph{and} the off-diagonal entry of $\Sigma$. The
analytical decomposition of the log-likelihood gain of VAR over
independent AR,
\begin{equation}\label{eq:varar_decomp}
  \log L_\text{VAR} - \log L_\text{indep}
  \;=\; \underbrace{\tfrac{T}{2}\log\!\tfrac{1}{1-\rho^2}}_{\text{from } \Sigma_{12}}
  \;+\; \text{(gain from cross-lags)},
\end{equation}
with $\rho$ the contemporaneous residual correlation, makes this
explicit. Consequently an AIC-vs-RMSE divergence is diagnostic: if
AIC prefers the bivariate model decisively but test-set RMSE is a
tie, the multivariate advantage lives in contemporaneous residual
correlation (a shared-forcing / entrainment signature) rather than
in cross-lag predictive structure (the coupling signature our paper
targets).
